% CREATED BY MAGNUS GUSTAVER, 2020
\chapter{Conclusion}
\label{chap:conclusion}
\label{chap:discussion}

This chapter closes the thesis by returning to the research questions introduced
in Chapter~\ref{chap:introduction}. Chapter~\ref{chap:experiment-analysis}
showed that \sys{} improves accuracy and returned-evidence compactness on
LoCoMo, LongMemEval, and \bench{}, with the largest relative gain on the
diagnostic temporal-bridge setting. The role of this chapter is not to repeat
those tables, but to explain what the results imply for long-term conversational
agent memory, where the current method remains limited, and what ethical and
deployment issues remain before such a system can be used as personal memory
infrastructure.

\section{RQ1}
\label{sec:conclusion-rq1}

\textbf{RQ1.} How can a long-term conversational memory system construct
evidence at query time while preserving the original memory structure and
avoiding heavy query-agnostic interpretation during ingestion?

\sys{} answers this question by treating memory access as query-time evidence
construction rather than write-time memory rewriting. The system does not rely
on a single precomputed memory record as the final representation of a past
conversation. Instead, it keeps the original session structure available, uses
lightweight indexing and temporal support at ingestion time, and delays the
main interpretation step until a concrete user query is available. This design
is important because the future question determines which part of a memory is
evidence. A detail that looks unimportant at write time may become the key
bridge for a later temporal or multi-hop question.

The Temporal Evidence Pool is the main mechanism that makes this answer
concrete. It stores compact events rather than raw sessions. Each event is tied
to a source session and carries both session time and event time. The memory
module therefore returns an inspectable evidence state to the downstream answer
model instead of a long, noisy history dump. This preserves source attribution
while reducing the burden on the answerer: the downstream model no longer has
to perform retrieval validation, evidence selection, temporal grounding,
compression, and final answer generation all at once.

The experimental results support this design. Compared with
\sys{}-Vanilla, \sys{}-GRPO improves accuracy while reducing the amount of
returned evidence on all three benchmarks. On LoCoMo, returned evidence drops
from 264 to 107 tokens; on LongMemEval, from 353 to 119 tokens; and on
\bench{}, from 287 to 234 tokens. These numbers should be interpreted as
returned-evidence compactness rather than full serving cost, but they show that
query-time evidence synthesis can improve answer usefulness without simply
passing more context downstream.

The answer to RQ1 is therefore that a long-term conversational memory system can
preserve the original memory structure by keeping ingestion lightweight and
moving the main interpretation into a query-time evidence provider. The memory
module should not be forced to choose, at write time, a single permanent
summary of what a session means. It can instead construct source-grounded,
temporally grounded evidence when the user query makes the relevant relation
visible.

\section{RQ2}
\label{sec:conclusion-rq2}

\textbf{RQ2.} How can a memory agent recover temporally grounded and multi-hop
evidence when the retrieval cue needed for the answer is not directly present
in the original query?

The central difficulty behind RQ2 is the evidence addressability gap. In many
long-term memory questions, the original query does not contain the phrase,
date, entity, or condition needed to retrieve the final answer-bearing session.
For example, a question may ask about what happened during the week of another
event. The system must first find the anchor event, derive a temporal or
semantic cue from it, and only then search for the destination evidence.

The system addresses this problem through an iterative
Planner--Retriever--Synthesizer loop. The important point is not merely that it
performs multiple retrieval rounds. The important point is that each round
updates the evidence state. Synthesized events from earlier rounds enter the
Temporal Evidence Pool, and the following round can use this pool as part of
the search state. This turns intermediate evidence into a new retrieval
address. When the first round finds a bridge event rather than the final
answer, the system has a structured way to carry that bridge forward.

Temporal grounding is essential to this process. Long-term conversational
memory has at least two relevant times: the time when a session occurred and
the time when the described event happened. These are often different. If a
user says on April 2 that a trip happened last week, the session time is April
2, but the event time is the previous week. \sys{} keeps this distinction in
the evidence state, allowing event time to guide later retrieval rather than
remaining only as unstructured text for the answer model.

The ablation results show that this evidence state is not an optional detail.
Removing the Temporal Evidence Pool lowers performance across the reported
tasks, with especially large drops on LongMemEval temporal reasoning and
\bench{}. Forcing single-round retrieval also hurts settings where the relevant
answer requires an intermediate bridge. At the same time, the results show that
iteration is not automatically beneficial for every question. Some questions
are answerable from first-round evidence, and extra search can introduce noise.
This is why the binary \texttt{sufficient}/\texttt{insufficient} verdict is a
core part of the design: the system must learn not only how to continue
searching, but also when to stop.

The answer to RQ2 is therefore that temporally grounded multi-hop evidence can
be recovered by maintaining a structured evidence state across retrieval rounds.
The memory agent needs to transform partial evidence into the next search cue,
while using event time and sufficiency judgments to avoid both premature
stopping and unnecessary over-search.

\section{RQ3}
\label{sec:conclusion-rq3}

\textbf{RQ3.} How can reinforcement learning be used to train the Synthesizer
to produce useful, compact, temporally grounded evidence and to judge whether
the current evidence is sufficient for downstream answering?

The role of reinforcement learning in \sys{} is to align the Synthesizer with
the memory-side responsibilities defined by the architecture. The Synthesizer
does not directly answer the user. It extracts and compresses evidence,
grounds events in source sessions and event times, and predicts whether the
current evidence is sufficient. GRPO is used to train this behavior inside the
same iterative environment in which the Synthesizer will be used at inference
time.

The reward design reflects this separation of responsibilities. The outcome
reward measures whether the accumulated evidence helps a frozen downstream
answerer produce the correct answer. The verdict reward trains the binary
sufficiency decision against gold evidence coverage. The length reward
encourages compactness when the answer outcome is
correct. The format gate enforces parseability at the trajectory level: if a
trajectory is malformed, it receives the fixed format penalty and no other
reward terms are added. This design avoids treating format as just another
additive preference.

The reward ablation supports the need for these components. Removing the format
and length rewards causes returned evidence to become much longer, showing that
compactness does not emerge reliably from answer accuracy alone. Removing the
verdict reward produces a different failure mode: some temporal categories can
remain strong, but multi-hop and hidden-bridge settings degrade, verdict
accuracy drops, and the system tends to search more while returning longer
evidence. This pattern is consistent with the task structure. When first-round
evidence is enough, outcome-only training can be effective; when the system
must decide whether a bridge is still missing, sufficiency supervision becomes
more important.

The answer to RQ3 is therefore that reinforcement learning is useful when it is
matched to the memory interface. GRPO does not replace the need for an
evidence-state architecture. Instead, it trains the Synthesizer to operate
within that architecture: produce parseable evidence, keep it compact, preserve
temporal grounding, and decide whether the accumulated evidence is enough for
downstream answering.

\section{Implications}
\label{sec:conclusion-implications}

The first implication is that long-term agent memory should not be evaluated
only by storage capacity or context length. A larger context window can expose
more history to the model, but it does not guarantee that the model will find
the right evidence, resolve the right time relation, or ignore irrelevant
sessions. The results on LongMemEval and \bench{} show that full-context
prompting can fail even when the answer-bearing information is present. The
hard part is not merely access to history, but construction of usable evidence.

The second implication is that the output boundary of a memory module matters.
If the memory module returns raw sessions, the answerer has to solve too many
subproblems at once. If the memory module directly returns a final answer, the
evidence construction process becomes hard to inspect. \sys{} takes a narrower
position: memory returns evidence and a verdict, while final answer generation
remains outside the memory module. This boundary makes errors easier to
diagnose. A failure can be attributed more specifically to retrieval recall,
temporal grounding, premature sufficiency, evidence bloat, or downstream answer
generation.

The third implication is that event time should be treated as search state, not
only as metadata. In conversational memory, time is often the thing that makes a
hidden dependency addressable. A past event may be mentioned after it occurred,
or an answer-bearing session may be reachable only through a time window
derived from another event. Keeping event time in the Temporal Evidence Pool
allows time to participate in retrieval, rather than appearing only in the
final prompt as text for the answerer to reason over.

The fourth implication is that query-time and write-time memory methods should
be seen as complementary rather than mutually exclusive. Write-time memory is
useful for stable user preferences, profiles, and frequently reused facts. The
results of this thesis suggest that it is less suitable as the only mechanism
for questions whose relevance depends on a future query. Query-time evidence
synthesis gives the system a way to interpret the memory in relation to the
current information need, while still preserving source grounding.

\section{Limitations and Future Work}
\label{sec:conclusion-limitations-future}

\subsection{Limitations}
\label{subsec:conclusion-limitations}

The first limitation is retrieval recall. \sys{} can only synthesize evidence
from sessions that are retrieved. If the Planner forms the wrong cue, the
retriever misses the answer-bearing session, or the candidate budget excludes a
relevant memory, the Temporal Evidence Pool cannot recover information that was
never observed. This is a general upper bound for search-based memory systems.

The second limitation is temporal normalization and event-time grounding.
Deterministic temporal normalization helps with common relative-time
expressions, but it does not cover all natural language time references, fuzzy
intervals, recurring events, time zones, cultural calendars, or changing future
plans. The Synthesizer can also make grounding errors by confusing session time
with event time or by compressing a broad interval into a narrow date. Because
event time is used as search state, these errors can propagate into later
retrieval rounds.

The third limitation is the expressiveness of the structured schema. The final
schema uses source-grounded events and a binary verdict. This makes parsing,
training, and evaluation more stable, but it gives limited room for
uncertainty, conflict, partial support, negative evidence, or multi-source
inference. The schema is therefore a pragmatic research interface, not a final
representation for all personal memory systems.

The fourth limitation is evaluation dependence. Final answer accuracy measures
whether the evidence helps a fixed downstream answerer produce the right
answer, but it is not a direct measure of evidence faithfulness. A correct
evidence state can still lead to a wrong answer, and a strong answerer can
sometimes compensate for incomplete evidence. The GPT-5.1 judge provides a
consistent automatic evaluation boundary, but any LLM judge may prefer certain
wordings, miss equivalent answers, or fail to detect source-grounding errors.

The fifth limitation is benchmark coverage. LoCoMo and LongMemEval are
important public long-term memory benchmarks, and \bench{} isolates the
temporal-bridge problem targeted by this thesis. However, these benchmarks do
not cover every realistic personal memory scenario. Real assistants must handle
user correction, memory deletion, stale preferences, cross-application data,
privacy policies, adversarial queries, and long-term drift. The present work
should therefore be read as a research prototype and diagnostic evaluation, not
as a complete deployed memory service.

The sixth limitation is cost. \sys{} reduces the evidence passed to the
downstream answerer, but the loop itself may require multiple retrieval rounds,
Synthesizer calls, and answerer or judge calls during evaluation. The token
compactness results in Chapter~\ref{chap:experiment-analysis} measure returned
evidence length, not full end-to-end latency, cost, or energy use.

\subsection{Future Work}
\label{subsec:conclusion-future-work}

One direction is to train the Planner and Retriever. This thesis focuses on
training the Synthesizer while keeping the rest of the environment fixed.
Future work could train a planner to generate better temporal and multi-hop
cues, or train a retrieval policy that dynamically allocates budget across
dense retrieval, lexical retrieval, temporal constraints, and source
constraints.

A second direction is to extend the evidence schema. Future systems could add
uncertainty, conflict, negative evidence, relation type, source confidence, and
multi-source support. Such fields should be added only when matching
supervision and evaluation are available; otherwise, the schema may become more
complex without becoming more reliable.

A third direction is stronger temporal reasoning. Rule-based normalization can
be expanded to more languages, time zones, fuzzy intervals, recurring events,
and calendar constraints. More robust variants could combine LLM-based
evidence synthesis with symbolic temporal reasoning or constraint checking, so
that generated event times are verified before they influence later retrieval.

A fourth direction is direct evidence evaluation. Future work should measure
source faithfulness, event-time correctness, evidence coverage, minimality,
verdict reliability, and uncertainty calibration directly. This would separate
retrieval success, evidence synthesis success, and downstream answer success
more cleanly than final answer accuracy alone.

A fifth direction is real longitudinal deployment evaluation. A realistic
personal memory system must support update, inspection, correction, deletion,
retention, and user control. Long-running user studies or high-fidelity
simulated agent environments are needed to test whether users can understand,
trust, and correct structured memory evidence over time.

A final direction is multilingual and cross-domain memory. The current
evaluation is centered on English benchmarks. Relative time, reference, and
calendar conventions differ across languages and domains. Testing Chinese,
Swedish, multilingual conversations, workplace agents, educational tutors, code
agents, and research assistants would clarify which parts of the evidence
addressability gap are general and which are domain-specific.

\section{Ethics}
\label{sec:conclusion-ethics}

Long-term conversational memory is privacy-sensitive by design. A system that
stores and searches user histories can accumulate information about health,
finances, relationships, work plans, location patterns, beliefs, and personal
vulnerabilities. The risk is not limited to storing individual sessions.
Query-time evidence synthesis can connect information across sessions and
reveal patterns that were never stated in one place.

Source-grounded evidence has both benefits and risks. On the positive side,
NanoMem makes memory use more auditable by returning source identifiers,
session times, and event times. This can help users and developers inspect why
an answer was produced. On the negative side, returning evidence can expose
private material if access control is weak or if the Synthesizer includes
irrelevant details. Compact evidence reduces exposure compared with raw history
injection, but it does not replace consent, access control, or policy-aware
retrieval.

Users should be able to control the memory lifecycle. A deployed system should
support consent, inspection, correction, deletion, retention limits, and clear
explanations of when memory is used. This is especially important because
personal memory can become stale. A preference, relationship, address, medical
detail, or work plan may change, and an agent that retrieves old evidence
without temporal qualification can give misleading or harmful advice.

High-risk domains require additional safeguards. In medical, legal, financial,
employment, or mental-health settings, memory evidence should be treated as
supporting context rather than authoritative truth. Premature sufficiency,
event-time drift, and source-attribution errors can cause real harm. Systems
using \sys{}-like evidence synthesis in such settings should require human
verification, domain-specific policies, and conservative refusal or escalation
behavior when evidence is incomplete.

Sustainability is also part of responsible deployment. \sys{} can reduce the
evidence tokens passed downstream, but this does not prove lower end-to-end
energy use. Training with GRPO, multi-round retrieval, LLM-based answering, and
automatic judging all consume computational resources. A deployed system should
therefore use budget-aware control, caching where appropriate, smaller or local
models for low-risk steps when possible, and transparent reporting of model
calls and evaluation costs.

AI tools were used during the preparation of this thesis for drafting support,
language editing, code-oriented checking, and consistency review. The author
remains responsible for the final text, claims, citations, experimental
descriptions, and conclusions. The reported experiments, tables, and method
descriptions were checked against the project artifacts and the completed
NanoMem research paper before inclusion in this thesis.

Overall, this thesis shows that long-term agent memory can be designed as
query-time evidence construction rather than raw history injection or opaque
answer generation. \sys{} makes progress toward compact, source-grounded, and
temporally grounded memory access, especially for temporal and multi-hop
queries. At the same time, reliable personal memory requires more than
benchmark accuracy. It also requires stronger retrieval, better evidence
verification, user control, privacy governance, and careful deployment in
settings where memory errors can affect real decisions.
